% !TeX root = ../sections/02_exercises.tex

\begin{exercise}
	R-8. 環 $A=\mathbb{Z}[\sqrt{3}]$ のイデアル
	\[
	I=(2\sqrt{3})
	\]
	は $2\sqrt{3}$ で生成された単項イデアルとする。
	\begin{enumerate}
		\item $\mathbb{Z}$-加群として
		\[
		I=6\mathbb{Z}\oplus 2\sqrt{3}\mathbb{Z}
		\]
		であることを示せ。
		\item $\mathbb{Z}$-加群として $A/I$ はねじれ加群で，
		\[
		A/I\cong \mathbb{Z}/2\mathbb{Z}\oplus \mathbb{Z}/6\mathbb{Z}
		\]
		に同型，すなわち不変系は $(2,6)$ であることを示せ。
	\end{enumerate}
\end{exercise}

\begin{background}
	$A=\mathbb{Z}[\sqrt{3}]$ は，
	$\mathbb{Z}$-加群として
	\[
	A=\mathbb{Z}\oplus \mathbb{Z}\sqrt{3}
	\]
	と表される。
	つまり，任意の元は一意的に
	\[
	a+b\sqrt{3}
	\qquad
	(a,b\in\mathbb{Z})
	\]
	と書ける。
	
	イデアル
	\[
	I=(2\sqrt{3})
	\]
	は
	\[
	I=\{(2\sqrt{3})x\mid x\in A\}
	\]
	で定義される。
	
	したがって，$x=a+b\sqrt{3}$ とおいて，
	\[
	(2\sqrt{3})(a+b\sqrt{3})
	\]
	を計算すれば，$I$ の $\mathbb{Z}$-加群としての形が分かる。
\end{background}

\begin{strategy}
	任意の元 $x\in A$ を
	\[
	x=a+b\sqrt{3}
	\qquad
	(a,b\in\mathbb{Z})
	\]
	と書く。
	すると
	\[
	(2\sqrt{3})(a+b\sqrt{3})
	=
	2a\sqrt{3}+6b
	\]
	である。
	
	したがって
	\[
	I
	=
	\{\,6b+2a\sqrt{3}\mid a,b\in\mathbb{Z}\,\}
	=
	6\mathbb{Z}\oplus 2\sqrt{3}\mathbb{Z}
	\]
	となる。
	
	次に，$A=\mathbb{Z}\oplus\mathbb{Z}\sqrt{3}$ の基底
	\[
	1,\sqrt{3}
	\]
	に関して，$I$ は
	\[
	6\mathbb{Z}\oplus 2\mathbb{Z}\sqrt{3}
	\]
	である。
	したがって商を取ると
	\[
	A/I
	\cong
	(\mathbb{Z}/6\mathbb{Z})
	\oplus
	(\mathbb{Z}/2\mathbb{Z})
	\]
	となる。
	
	不変系は小さい方が前に来るように並べて
	\[
	(2,6)
	\]
	である。
\end{strategy}